% =====================================================================
%  config/packages.tex
%  Semua paket yang dipakai SELURUH buku didefinisikan di sini.
%  Bab tidak boleh memuat \usepackage sendiri (Aturan AI #3).
%  Isi file ini dipindahkan dari config/preamble.tex lama.
% =====================================================================

% --- Bahasa & pengkodean -------------------------------------------------
\usepackage[utf8]{inputenc}
\usepackage[T1]{fontenc}
% Latin Modern: font Type1 yang scalable. Diperlukan oleh font expansion
% microtype (font T1/EC bawaan bersifat bitmap pada sebagian ukuran).
\usepackage{lmodern}
\usepackage[indonesian]{babel}

% --- Matematika ----------------------------------------------------------
\usepackage{amsmath, amssymb, amsthm, mathtools}

% --- Grafik & gambar -----------------------------------------------------
\usepackage{graphicx}
\usepackage{float}   % menyediakan opsi [H] pada figure dan table
\usepackage{tikz}
\usetikzlibrary{arrows.meta, positioning, shapes.geometric, calc, fit, backgrounds}

% --- Tabel ---------------------------------------------------------------
\usepackage{booktabs}
\usepackage{array}
\usepackage{longtable}
\usepackage{tabularx}

% --- Warna & kode program ------------------------------------------------
\usepackage{xcolor}
\usepackage{listings}

% --- Daftar & kotak ------------------------------------------------------
\usepackage{enumitem}
\usepackage{tcolorbox}

% --- Kompilasi per bab ---------------------------------------------------
\usepackage{subfiles}

% --- Sitasi & daftar pustaka ---------------------------------------------
% natbib dimuat SEBELUM hyperref (Arsitektur_Buku_Ajar.md Bagian 20).
\usepackage{natbib}

% --- Tata letak halaman --------------------------------------------------
\usepackage{geometry}

% --- Tautan & PDF --------------------------------------------------------
\usepackage[hidelinks,breaklinks]{hyperref}
\usepackage{url}
% xurl memperbolehkan pemenggalan URL di sembarang karakter, sehingga URL
% panjang pada Lampiran C tidak menghasilkan underfull hbox. Dimuat setelah
% hyperref dan url.
\usepackage{xurl}

% --- Mikro-tipografi -----------------------------------------------------
% Dimuat setelah paket font. Character protrusion + font expansion sangat
% menekan jumlah overfull/underfull box tanpa mengubah isi.
\usepackage{microtype}

% --- Hyphenation (per bahasa) --------------------------------------------
% Pola hyphenation Indonesia (babel) tidak mencakup istilah Inggris/teknis
% berikut. Harus diberikan per-bahasa: \hyphenation biasa akan terdaftar
% untuk bahasa default, bukan untuk bahasa aktif (indonesian).
\babelhyphenation[indonesian]{su-per-in-creas-ing non-su-per-in-creas-ing Po-ly-al-pha-be-tic}

% --- Jaring pengaman overfull --------------------------------------------
% Memberi TeX kesempatan terakhir menambah rentangan antar-kata, tetapi
% hanya untuk paragraf yang jika tidak akan melampaui margin.
\setlength{\emergencystretch}{2em}
